\chapter{Conclusiones}\label{cap:conclusiones}

Tras finalizar el desarrollo de Deckbuilder hemos logrado cumplir el objetivo principal que era diseñar e implementar una aplicación funcional que utiliza las tecnologías de C\# y MongoDB.

Una de las características más importantes del proyecto ha sido el uso de una base de datos documental. Gracias al uso de MongoDB hemos podido gestionar de la mejor forma la variabilidad de los datos entre los distintos juegos de cartas que hemos usado, evitando de esta forma las limitaciones de los modelos relaciones que conocemos.

Viéndolo desde un punto de vista técnico, el desarrollo de un CRUD completo, junto con el frontend interactivo, nos ha permitido cubrir todo el flujo del usuario dentro de la aplicación, desde la autenticación hasta la gestión completa de las barajas y colecciones.

En conclusión, el proyecto nos ha permitido aplicar de forma práctica los conceptos estudiados en la asignatura, especialmente en el uso de bases de datos NoSQL, consolidando de esta forma una buena base con estas nuevas tecnologías.
